\chapter{CONSIDERAÇÕES FINAIS}

Este trabalho realizou uma análise comparativa entre a execução nativa de código Rust e sua execução compilada para WebAssembly, avaliando os impactos em tempo de execução e consumo de memória. A utilização de CPU também foi coletada, mas mostrou-se pouco informativa para a comparação, pelos motivos descritos na Seção~\ref{sec:consumo-memoria}, sendo utilizada apenas para identificar os períodos de execução.

Para embasar essa análise, foram abordados os fundamentos da linguagem Rust e as características do WebAssembly na revisão da literatura. Além disso, na etapa prática foram selecionados, do The Computer Language Benchmarks Game, 5 algoritmos de diferentes classes de problemas, implementados em Rust, portados para WebAssembly e executados em ambiente controlado via Docker em 3 máquinas distintas. Os resultados mostraram que o Rust nativo foi mais rápido que o WebAssembly em todos os algoritmos avaliados, com \emph{overhead} do WebAssembly entre 10,9\%, no fannkuch-redux, e 126,9\%, no mandelbrot, e mediana de 21,2\%. A análise de memória evidenciou que parte do consumo do ambiente WebAssembly não está associada à execução dos algoritmos, mas ao ambiente de execução, que mantém componentes como o Chrome Headless e o Puppeteer em funcionamento contínuo e eleva o consumo de memória em repouso para cerca de 210 MiB.

Diante disso, a utilização das métricas de \emph{baseline} e consumo líquido mostrou-se fundamental para isolar o impacto real dos algoritmos, permitindo uma análise mais precisa e justa entre os ambientes. Por essa métrica, o consumo líquido dos dois ambientes foi reduzido na maioria dos algoritmos, com o WebAssembly ligeiramente acima do Rust nativo, e o maior consumo em ambos foi o do mandelbrot, determinado pelo tamanho do \emph{buffer} de saída. Observou-se ainda que o \emph{overhead} de tempo varia conforme o padrão de acesso à memória, indicando que algoritmos com acesso sequencial e previsível tendem a apresentar desempenho mais próximo do nativo, enquanto aqueles com acesso irregular ou escrita intensiva apresentam maior degradação de desempenho.

Conclui-se, portanto, que o WebAssembly representa uma alternativa viável à execução nativa para aplicações Rust executadas no navegador com acesso predominantemente sequencial à memória, cenário em que o custo adicional ficou entre 10,9\% e 21,2\%. Para cargas com escrita intensiva ou acesso irregular à memória, o custo pode superar o dobro do tempo nativo, e a execução nativa segue mais indicada quando o desempenho é o principal requisito.

Como limitação, os algoritmos foram executados com uma única \emph{thread} nos dois ambientes, pois não foi possível habilitar o processamento paralelo no WebAssembly. Como trabalhos futuros, sugere-se avaliar a execução multi-\emph{thread} no navegador, ampliar o conjunto de algoritmos e comparar outros mecanismos de execução de WebAssembly.
